% !TEX program = lualatex
% JEFE 5 — Unicode y trucos de TeX: español, griego, ruso y alemán en el mismo documento, fuentes
% OpenType, matemática Unicode, química, programación en expl3 y en Lua, comandos con argumentos
% raros (xparse), catcodes cambiados, \csname, un archivo que se escribe al compilar (filecontents)
% y otro que se lee, marcadores del PDF con tildes y biber con nombres no latinos.
\begin{filecontents*}[overwrite]{constantes.tex}
\newcommand{\cLuz}{299\,792\,458}
\newcommand{\cPlanck}{6{,}626\,070\,15\times10^{-34}}
\end{filecontents*}
\begin{filecontents*}[overwrite]{fuentes.bib}
@book{landau, author = {Ландау, Лев Давидович and Лифшиц, Евгений Михайлович}, title = {Теоретическая физика},
  volume = {2}, publisher = {Наука}, location = {Москва}, year = {1988}}
@book{euclides, author = {Εὐκλείδης}, title = {Στοιχεῖα}, year = {-300}}
@book{gauss, author = {Gauß, Carl Friedrich}, title = {Disquisitiones Arithmeticae}, location = {Leipzig}, year = {1801}}
@article{nuñez, author = {Núñez Ñahui, José Ángel}, title = {Cálculo de señales en el año 2026}, journaltitle = {Revista Peruana de Ingeniería}, volume = {3}, year = {2026}}
\end{filecontents*}
\documentclass[11pt,a4paper]{article}
\usepackage[margin=2.2cm]{geometry}
\usepackage[ngerman,main=spanish,es-noshorthands]{babel}
\babelprovide[import, onchar=ids fonts]{greek}
\babelprovide[import, onchar=ids fonts]{russian}
\usepackage{fontspec}
\setmainfont{FreeSerif}
\setsansfont{FreeSans}
\setmonofont{FreeMono}
\usepackage{unicode-math}
\setmathfont{Latin Modern Math}
\usepackage{xparse}
\usepackage{chemformula}
\usepackage{siunitx}
\sisetup{output-decimal-marker={,}}
\usepackage{luacode}
\usepackage[backend=biber,style=alphabetic]{biblatex}
\addbibresource{fuentes.bib}
\usepackage[unicode]{hyperref}
\input{constantes}
\begin{luacode*}
  function fibonacci(n)
    local a, b, t = 0, 1, {}
    for i = 1, n do t[i] = a; a, b = b, a + b end
    tex.sprint(table.concat(t, ", "))
  end
\end{luacode*}

% ── expl3: una función que repite, una suma con l3fp y una lista ordenada ──
\ExplSyntaxOn
\NewDocumentCommand{\Repetir}{ O{,~} m m }{ \int_step_inline:nnn {1}{#2}{ ##1:~#3 \int_compare:nNnT {##1} < {#2} {#1} } }
\NewDocumentCommand{\Calcular}{ m }{ \fp_eval:n { round(#1, 6) } }
\clist_new:N \l_mis_valores_clist
\NewDocumentCommand{\Ordenar}{ m }{
  \clist_set:Nn \l_mis_valores_clist {#1}
  \clist_sort:Nn \l_mis_valores_clist { \int_compare:nNnTF {##1} > {##2} \sort_return_swapped: \sort_return_same: }
  \clist_use:Nn \l_mis_valores_clist {~<~}
}
\ExplSyntaxOff

% ── xparse: argumento opcional, estrella y delimitadores propios ──
\NewDocumentCommand{\Vector}{ s O{v} r() }{\IfBooleanTF{#1}{\boldsymbol{#2}}{\vec{#2}} = (#3)}

% ── \csname: nombres de comandos fabricados con un bucle ──
\makeatletter
\@for\x:=uno,dos,tres\do{\expandafter\xdef\csname numero\x\endcsname{\x}}
\makeatother

% ── catcodes: | como atajo de \texttt mientras dure el grupo ──
\newcommand{\activarbarra}{\catcode`\|=\active}
{\activarbarra\gdef|#1|{\texttt{#1}}}

\begin{document}
\title{Unicode y trucos de \TeX}
\author{José Ángel Núñez Ñahui \and Αλέξανδρος \and Дмитрий}
\date{\today}
\maketitle
\tableofcontents

\section{Varios idiomas en la misma página: ñ, ü, ß, α, ж}
Español con tildes y eñes: «¿Por qué el pingüino no vuela?», año, señal.

\foreignlanguage{ngerman}{Größe, Straße, Übergang — „Anführungszeichen“.}

Griego sin cambiar nada: η εξίσωση του Σρέντινγκερ είναι θεμελιώδης.

Ruso: уравнение Шрёдингера — основа квантовой механики~\cite{landau}.

Fórmulas en Unicode directo: $∇ × 𝐄 = -\frac{∂𝐁}{∂t}$, $∮_C 𝐁·d𝐥 = μ_0 I$, $c = \cLuz~\mathrm{m/s}$,
$h = \cPlanck~\mathrm{J·s}$ (leídos de un archivo escrito al compilar).

\section{Química y unidades}
\ch{2 H2 + O2 -> 2 H2O} \quad \ch{CuSO4 * 5 H2O} \quad \ch{Fe^{3+} + 3 OH- -> Fe(OH)3 v}
\quad la constante es \SI{8.314}{\joule\per\mole\per\kelvin}.

\section{Programación dentro del documento}
\begin{itemize}
  \item expl3 repite: \Repetir{3}{hola}. Con otro separador: \Repetir[ | ]{4}{x}.
  \item expl3 calcula: $\sqrt{2}+\pi = \Calcular{sqrt(2)+pi}$, $e^{\pi} - \pi = \Calcular{exp(pi)-pi}$.
  \item expl3 ordena: \Ordenar{42,7,19,3,88,1}.
  \item xparse: $\Vector(1,2,3)$ y con estrella $\Vector*[F](0,-9{,}8,0)$.
  \item \textbackslash csname: \numerouno, \numerodos, \numerotres.
  \item Lua: los 10 primeros números de Fibonacci son \directlua{fibonacci(10)}.
  \item catcode: {\activarbarra Esto |se ve como código| sin escribir el comando.}
\end{itemize}

\begin{luacode*}
  local filas = {}
  for n = 1, 8 do
    filas[#filas+1] = string.format("%d & %d & %.4f & %s \\\\", n, n*n, math.sqrt(n), (n % 2 == 0) and "par" or "impar")
  end
  tex.sprint("\\begin{center}\\begin{tabular}{rrrl}\\hline $n$ & $n^2$ & $\\sqrt n$ & paridad\\\\\\hline")
  tex.sprint(table.concat(filas))
  tex.sprint("\\hline\\end{tabular}\\end{center}")
\end{luacode*}

\section{Bibliografía con nombres en otros alfabetos}
Euclides~\cite{euclides}, Gauß~\cite{gauss}, Landau~\cite{landau} y un autor peruano~\cite{nuñez}.
\printbibliography
\end{document}
